\chapter{The system as built}
\label{ch:system}

This chapter describes what exists, at the level of what each part is for.
The two chapters after it go into the parts that carry the most risk.

\section{The hardware}

Two Kinova Gen3 arms, each with seven joints, a two-finger gripper and a small
colour-and-depth camera at the wrist, are carried on a backpack frame. The
mounting brackets sit behind the wearer's shoulders and tilt the arms forward
so they come over the top rather than out to the sides.

The driver's mannequin arm is a scale model of the same joint sequence with a
rotary sensor at each joint and a small inertial sensor at the wrist. A Teensy
4.1 microcontroller~\cite{pjrc2020teensy41} reads all of it and streams the
numbers over a USB cable at about \SI{50}{\hertz}. There is no force feedback
of any kind: the driver feels the mannequin arm's own friction and nothing
about what the robot is touching.

Everything runs on one laptop under Linux, with the robot software built on
ROS~2~\cite{macenski2022ros2}. That machine has a small graphics card with
\SI{4}{\giga\byte} of memory, which decided several questions later on by
ruling out the larger vision and planning packages outright.

\begin{figure}[htbp]
  \centering
  \photo{58mm}{The platform as assembled: the backpack frame worn on a stand,
  both arms mounted and folded to the presentation pose, the mannequin master
  on the bench beside it with its cable run, and the emergency stop within
  reach. One wide shot, and one close view of a shoulder bracket showing the
  \SI{60}{\degree} forward tilt.}
  \caption[The platform]{The platform. A photograph is needed here rather than
  a rendering, because the mounting bracket and the cable routing are the two
  things a reader cannot infer from the model.}
  \label{fig:platform}
\end{figure}

\section{The software, in one picture}

The software is about forty small programs that talk to each other over named
channels. They are grouped into eight packages by what they are responsible
for, and the grouping matters for one reason: the teleoperation package is not
allowed to depend on any of the others.

That rule sounds like tidiness and is not. Plain teleoperation is the baseline
that every assisted condition in the study is compared against. If the
teleoperation code called into the autonomy code, the comparison would be
between the autonomy and itself.

\begin{table}[htbp]
  \centering
  \small
  \caption[Software packages]{The eight first-party packages and what each is
  responsible for. Sizes are measured lines of Python.}
  \label{tab:packages}
  \begin{tabular}{@{}p{0.20\linewidth}p{0.60\linewidth}r@{}}
    \toprule
    \textbf{Package} & \textbf{Responsible for} & \textbf{Lines} \\
    \midrule
    Teleoperation & Reading the master arm, turning it into commands, the
      clutch, motion smoothing, the emergency stop, the operating window &
      25\,737 \\
    Perception & Cameras, finding objects, measuring the table, building a map
      of the scene, tracking the wearer & 9\,431 \\
    Autonomy & Choosing a grasp, guessing the goal, handing control over,
      understanding a sentence & 3\,886 \\
    Experiments & Tasks, conditions, logging, sessions, analysis & 12\,278 \\
    Description & The model of the arms, the frame and the wearer & --- \\
    Motion planning & Inverse kinematics, controllers, collision settings &
      --- \\
    Controller bridge & The hand controllers and their safety checks &
      2\,709 \\
    Controller autonomy & Assistance on the controller path & 246 \\
    \bottomrule
  \end{tabular}
\end{table}

\section{From a hand movement to a moving arm}

\Cref{fig:datapath} is the path every command takes, whichever way the arms
are being driven. It is worth reading once carefully, because most of the rest
of this thesis is about one or another of its boxes.

\begin{figure}[htbp]
  \centering
  % From the driver's hand to the arm, with every gate that can stop it.
\begin{tikzpicture}[x=1mm, y=1mm, font=\scriptsize, node distance=6mm]

  \node[sense] (in) at (0,0) {Driver's input\\ \tiny(mannequin arm, hand
        controllers, or a typed sentence)};
  \node[stage, below=7mm of in] (pose) {Wanted hand position and direction};
  \node[choice, below=7mm of pose, text width=26mm] (ik) {Can the arm hold
        that pose?};
  \node[guard, below=9mm of ik] (clear) {How close does the whole arm come to
        the wearer?};
  \node[stage, below=7mm of clear] (gen) {Smooth the motion:\\ speed, acceleration
        and jerk limited};
  \node[motor, below=7mm of gen] (out) {Joint commands to the arms};

  \draw[go] (in) -- (pose);
  \draw[go] (pose) -- (ik);
  \draw[go] (ik) -- node[tag, right]{yes} (clear);
  \draw[go] (clear) -- node[tag, right]{above \SI{150}{\milli\metre}} (gen);
  \draw[go] (gen) -- (out);

  \node[guard, right=26mm of ik, text width=34mm] (retry)
        {Try again with the elbow swung to a different position};
  \draw[halt] (ik.east) -- node[stoptag, above]{no} (retry.west);
  \draw[go] (retry.north) |- (pose.east);

  \node[guard, right=26mm of clear, text width=34mm] (hold)
        {Hold. Publish nothing. Say which body part was too close, by name};
  \draw[halt] (clear.east) -- node[stoptag, above]{below} (hold.west);

  \node[guard, left=24mm of gen, text width=32mm] (estop)
        {Emergency stop, dead-man, or any of nine named blocking conditions};
  \draw[halt] (estop.east) -- (gen.west);

  \node[person, left=24mm of clear, text width=32mm] (body)
        {Where the wearer's body is, from the camera when it is trusted and
        from the stored model otherwise};
  \draw[go, violet!60!black] (body.east) -- (clear.west);

\end{tikzpicture}

  \caption[The command path]{The path from the driver's input to the arms.
  Red boxes can stop the motion; the green box is the only one that moves it.
  Two of the red boxes matter most: the retry that swings the elbow to a
  different position when the first attempt fails, and the clearance check,
  which is measured against the wearer's body geometrically rather than being
  taken from the collision settings.}
  \label{fig:datapath}
\end{figure}

The path is deliberately short. The wanted hand pose goes to a solver that
returns joint angles, the whole arm is checked against the wearer, the result
is smoothed so no joint is asked to move faster than it can, and the joint
commands go out. There is no planner in that loop and there should not be: the
driver's hand is the plan, and the right problem is to follow it smoothly. For
the autonomous mode, where there is a goal to plan towards, a planner was
added and is described in \cref{ch:perception}.

\section{Three ways to drive}

\begin{figure}[htbp]
  \centering
  % The three ways of driving, and what each supplies.
\begin{tikzpicture}[x=1mm, y=1mm, font=\scriptsize]

  \foreach \i/\name/\input/\assist in {
     0/{Mannequin master}/{driver moves a sensed arm by hand}/{none: the driver supplies everything},
     1/{Hand controllers}/{two tracked controllers held across the room}/{none},
     2/{A typed or spoken sentence}/{``put the blue block on the blue mat''}/{the machine plans and executes}}
  {
    \pgfmathsetmacro{\x}{\i*54}
    \node[sense, minimum height=13mm, text width=44mm] at (\x, 0)
      (a\i) {\textbf{\name}\\[2pt] \input};
    \node[stage, minimum height=13mm, text width=44mm, below=6mm of a\i]
      (b\i) {\assist};
  }

  \node[stage, minimum height=9mm, text width=112mm] (common) at (54,-38)
    {\textbf{One common path from here on.} The same wanted hand pose, the
     same inverse kinematics, the same clearance check, the same motion
     smoothing, the same emergency stop.};

  \foreach \i in {0,1,2} { \draw[go] (b\i) -- (common); }

  \node[motor, minimum height=8mm, text width=60mm] (arm) at (54,-52)
    {The arms};
  \draw[go] (common) -- (arm);

  \node[tag, text width=110mm, align=left] at (54,-62)
    {Assistance is added only above the common path, never inside it. Whichever
     way the arms are being driven, the part that keeps the wearer safe is the
     same code running the same checks.};

\end{tikzpicture}

  \caption[The three ways of driving]{The three input routes. They differ only
  above the common path. Everything that protects the wearer is shared, so a
  change to the safety layer cannot apply to one condition of the study and
  not another.}
  \label{fig:modes}
\end{figure}

The mannequin master is the baseline. The driver moves it and the robot copies
the movement of its tip, scaled and offset.

The hand controllers came from a virtual reality headset, but the way they are
used has almost nothing to do with virtual reality. Nobody wears the headset.
It stands on a shelf and acts as the tracking reference while the driver
holds the two controllers and watches the real robot across the room. That
arrangement is the one the laboratory actually uses, and it caused two
specific safety problems that are dealt with in \cref{sec:tracking}.

The third route takes a sentence. Typing \emph{put the blue block on the blue
mat} produces a plan, shows the driver what it understood before anything
moves, and waits for a confirming click. It is built on a grammar rather than
a language model, for three reasons: a grammar can explain why it refused, it
can be tested exhaustively against a list of phrasings with a scored outcome
for each, and it fits in the memory the graphics card has left after the
object detector has loaded.

\section{The operating window}

Everything reachable from a terminal is also reachable from one window, and
that is a rule rather than a convenience. A feature that can only be started
by typing a command does not exist for the person running a session, who has
a wearer standing in a rig and no attention to spare.

\begin{figure}[htbp]
  \centering
  \screenshot{78mm}{The main operating window, full screen, with a
  simulation running. It should show: the left column grouped into SET UP,
  RUN and STATUS; the embedded three-dimensional view showing the commanded
  arms; the panel beside it drawn from the arms' reported positions; and the
  divergence readout underneath both. Take the shot with both arms away from
  the home pose so the two views visibly differ.}
  \caption[The operating window]{The window used to run a session. The point
  of the side-by-side views is that the left one shows what the arms were
  told to do and the right shows what they report doing, so a difference is
  visible without reading a log.}
  \label{fig:gui}
\end{figure}

\begin{figure}[htbp]
  \centering
  \screenshot{62mm}{The simulated scene as the operator sees it: the wearer's
  body drawn as the shapes the clearance check actually uses, the table, the
  four coloured blocks and the two coloured mats, the marked working area, and
  both arms mid-task. A second panel showing the same scene with the
  clearance shell drawn around the wearer would make the margin visible.}
  \caption[The simulated scene]{The simulation. The wearer is drawn from the
  same description the safety check measures against, so what is on screen and
  what is being enforced cannot drift apart.}
  \label{fig:sim}
\end{figure}

The window is checked by a program that presses every control in it and then
checks what the press did, currently 257 separate checks. That sounds
excessive until one notices that the first time it ran end to end it found
four faults, including a set of buttons that had been doing nothing at all
because the error they raised was being swallowed silently.
